\section{Background}


\subsection{Threshold Signatures}

\begin{definition}[$(t, n)$-Threshold Signature Scheme]
A \TSS{} distributes a signing key $sk$ among $n$ parties such that any $t$ parties can collaboratively produce a valid signature, but fewer than $t$ parties learn nothing about $sk$.
\end{definition}

\subsection{CGGMP21}

The Canetti-Gennaro-Goldfeder-Makriyannis-Peled protocol~\cite{cggmp2021} provides threshold ECDSA signing with:

\begin{itemize}
    \item \textbf{UC security}: Composably secure in the universal composability framework
    \item \textbf{Identifiable abort}: If signing fails, the malicious party is identified
    \item \textbf{Pre-signing}: Expensive operations done offline; online signing is fast
\end{itemize}

Signing complexity: $O(t^2)$ messages in the online phase, $O(n \cdot t)$ in pre-signing.

\subsection{FROST}

Flexible Round-Optimized Schnorr Threshold~\cite{komlo2020frost} provides threshold Schnorr/EdDSA signing with:

\begin{itemize}
    \item \textbf{Two rounds}: Commitment round + signing round
    \item \textbf{No trusted dealer}: Distributed key generation via Pedersen DKG
    \item \textbf{Robustness}: Optional identifiable abort extension
\end{itemize}

FROST is 3--5$\times$ faster than CGGMP21 for equivalent threshold parameters because Schnorr signatures have simpler algebraic structure than ECDSA.

\subsection{Lux M-Chain}

The Lux Network's M-Chain~\cite{lux-mchain-mpc} provides native MPC infrastructure:

\begin{itemize}
    \item Key generation ceremonies coordinated on-chain
    \item Threshold signing requests submitted as transactions
    \item Key shares stored in TEE-protected enclaves on validator nodes
    \item Cross-chain signing: M-Chain key shares can sign for any Lux chain
\end{itemize}

Zoo Network inherits M-Chain access as an L2, enabling threshold signing without running separate MPC infrastructure.

